\lecture{9}
\title{Gradient Descent}

\begin{document}

\subsection{Gradient Descent}

Recall from \href{/6.790/lectures/8/#mle-logistic-regression}{Lecture 8}
that logistic regression requires minimizing
the following (logistic) loss:
\[ \mathcal{L}(w) = \sum_{i = 1}^n \log_2\left[
    1 + \exp\left(-y^{(i)}w^{\top}x^{(i)}\right)
\right]. \]
Unfortunately, unlike with linear regression,
there is no closed-form expression
for the minimizing parameter of logistic loss.
But because logistic loss is \emph{convex},
we can apply \textbf{gradient descent} to find the global minimum.
(Recall from \href{/6.790/lectures/8/#addendum-convexity}{Lecture 8}
that we can say ``\emph{global} minimum'' rather than
``\emph{local} minimum'' because of convexity.)

\textbf{Definition.}
Let $f \colon \mathbb{R}^n \to \mathbb{R}$ be a differentiable function.
In \textbf{gradient descent (GD)},
we start with some $w_0 \in \mathbb{R}^n$
and proceed via:
\[ w_{t + 1} \leftarrow w_t - \eta \nabla f(w_t). \]
Note the hyperparameter $\eta > 0$,
which we call the \textbf{step size}.
The hope is that $\{w_t\}_{t = 0}^{\infty}$
converges to a local minimizer of $f$.

\subsection{Stationary Points}

Gradient descent stops at a value $w^* \in \mathbb{R}^n$
where $\nabla f(w^*) = 0$; we call such a point $w^*$
a \textbf{stationary point}.
We can think of stationary points as
being classifiable into two categories:
\begin{itemize}
    \item Category I: Global maximizers, global minimizers,
    local maximizers, or local minimizers.
    \item Category II: Saddle points
    (as in, stationary points that are not Category I).
\end{itemize}
This is kind of irrelevant, though;
gradient descent only converges to saddle points
for a measure-zero set of initial points $w_0$.

For a further discussion of how the Hessian relates
to saddle points and whatnot, see
\href{/18.701/lectures/16/#foreshadowing-quadratic-polynomials-are-bilinear-forms}{Lecture 16}
and \href{/18.701/lectures/19/#spectral-application-1-classifying-quadric-surfaces}{Lecture 19}
of 18.701 (specifically the subsections linked above).

\subsection{Lower-Bounding GD Efficiency}

For one-dimensional functions $f \colon \mathbb{R} \to \mathbb{R}$,
we can upper-bound the number of steps required
to get close to a stationary point.

\begin{theorem}{One-Dimensional Descent Lemma}
    Suppose $f \colon \mathbb{R} \to \mathbb{R}$ is twice differentiable,
    and $|f''(w)| \leq \beta$ for all $w$.
    Then if we set $\eta = \frac{1}{\beta}$,
    the iterates of gradient descent satisfy:
    \[ f(w_{t + 1}) - f(w_t) \leq -\frac{1}{2\beta} f'(w_t)^2. \]
    In particular, each step of gradient descent
    produces an improvement
    of at least $\frac{1}{2\beta} f'(w_t)^2$.
\end{theorem}

\begin{proof}
    We won't provide a rigorous proof.
    The rough idea is to use the quadratic Taylor approximation.
    \[
        f(w_{t + 1})
        \approx f(w_t) + \delta f'(w_t) + \frac{\delta^2}{2}f''(w_t)
        ~ \text{ where } \delta = -\eta f'(w_t).
    \]
    Plugging in $\delta = -\eta f'(w_t)$
    and rearranging yields:
    \[ f(w_{t + 1}) - f(w_t)
    \approx - \eta f'(w_t)^2 \left[ 1 - \frac{\eta}{2} f''(w_t) \right]
    \leq -\frac{1}{2\beta} f'(w_t)^2
    ~~~ \text{ assuming }
    \eta = \frac{1}{\beta}
    \text{ and }
    |f''(w)| \leq \beta. \]
    That's the promised bound, with many computations omitted.
    This proof, of course, lacks rigor in that
    everything carries an $\approx$ symbol
    due to the Taylor approximation. ~ $\blacksquare$
\end{proof}

An immediate corollary of the Descent Lemma is the following.

\begin{theorem}{One-Dimensional GD Efficiency}
    Suppose, in the setting of the Descent Lemma,
    that $f(w_0) - \min_w f(w) = M$.
    If gradient descent yields $\{w_t\}_{t = 0}^{\infty}$,
    then for any $\epsilon > 0$,
    there is some $T \leq \frac{2 \beta M}{\epsilon^2}$
    for which $|f'(w_T)| < \epsilon$.
\end{theorem}

\begin{remark}
    Note that none of the reasoning here
    carries an assumption that $f$ is convex.
\end{remark}

For some reason, 6.790 does not state
the multivariate versions of these theorems,
even though their statements and proofs
are nearly identical to the one-dimensional case.
We state them here for completeness.

\begin{theorem}{Multivariate Descent Lemma}
    Suppose $f \colon \mathbb{R}^n \to \mathbb{R}$ is twice differentiable,
    and $-\beta I \preceq \nabla^2 f(w) \preceq \beta I$ for all $w$.
    Then if we set $\eta = \frac{1}{\beta}$,
    the iterates of gradient descent satisfy:
    \[ f(w_{t + 1}) - f(w_t) \leq -\frac{1}{2\beta} \|\nabla f(w_t)\|^2. \]
    In particular, each step of gradient descent
    produces an improvement
    of at least $\frac{1}{2\beta} \|\nabla f(w_t)\|^2$.
\end{theorem}

\begin{remark}
    The ordering $A \preceq B$ is
    the \emph{Loewner order} on symmetric matrices:
    $A \preceq B$ if and only if $B - A$ is positive semidefinite.
\end{remark}

\begin{theorem}{Multivariate GD Efficiency}
    Suppose, in the setting of the Descent Lemma,
    that $f(w_0) - \min_w f(w) = M$.
    If gradient descent yields $\{w_t\}_{t = 0}^{\infty}$,
    then for any $\epsilon > 0$,
    there is some $T \leq \frac{2 \beta M}{\epsilon^2}$
    for which $\|\nabla f(w_T)\| < \epsilon$.
\end{theorem}

\subsection{Linear Regression Convergence Rate}

Suppose we use gradient descent
to find the minimum in simple least-squares linear regression.
Then an overly large step size
can cause gradient descent to fail to converge!

\begin{center}
    \frame{\includegraphics[width=0.6\textwidth]{lecture-09/GD-failure}}
\end{center}

This often happens when the level sets of $f$
are elliptically-shaped, rather than spherically-shaped.

\begin{center}
    \frame{\includegraphics[width=0.4\textwidth]{lecture-09/level-sets}}
\end{center}

Let's investigate how this behavior appears
specifically in the case of least-squares linear regression.

\begin{theorem}{Least-Squares Gradient Descent}
    Suppose we use gradient descent to find
    a $w^* \in \mathbb{R}^d$ such that,
    given training data $X \in \mathbb{R}^{n \times d}$
    and $y \in \mathbb{R}^n$,
    the loss $f(w) := \|Xw - y\|^2$ is minimized.
    Then the sequence $\{w_t\}_{t = 0}^{\infty}$ satisfies:
    \[ w_t - w^* = (I_d - \eta H)^t (w_0 - w^*), \]
    where $H := \nabla^2 f(w) = 2X^{\top}X$
    is the Hessian matrix (notably, independent of $w$).
\end{theorem}

\begin{proof}
    It's just plain computation.
    First, compute the relevant derivatives.
    \[ \nabla f(w) = 2X^{\top}(Xw - y)
    ~~~ \text{ and } ~~~
    \nabla^2 f(w) = 2X^{\top}X. \]
    Then note the minimizer $w^*$ is known to satisfy
    (see \href{/18.650/lectures/23/#least-squares-loss}{Lecture 23}
    of 18.650) the equation $X^{\top}Xw^* = X^{\top}y$, so:
    \[ w_{t + 1} - w^*
    = \left[w_t - \underbrace{2 \eta \left(
        X^{\top}X w_t - X^{\top}y
    \right)}_{
        = \ \eta \nabla f(w_t)
    }\right] - w^*
    = \left(I_d - 2 \eta X^{\top}X\right)(w_t - w^*). \]
    Recalling $H = 2X^{\top}X$ and applying induction
    finishes the job. ~ $\blacksquare$
\end{proof}

For very large $t$, the effect of $(I_d - \eta H)^t$
can be summarized by its eigenvalues.
Say $\lambda_1 \geq \dots \geq \lambda_d > 0$
are the eigenvalues of $H$.
Then we can make the following observations.
\begin{enumerate}
    \item The eigenvalues of $I_d - \eta H$
    are precisely $\left\{1 - \eta \lambda_i\right\}_{i = 1}^d$.
    \item To ensure convergence,
    we require all eigenvalues to be in $(-1, 1)$.
    \begin{itemize}
        \item The \emph{largest} eigenvalue $\lambda_1$
        gives the tightest bound here;
        the constraint $1 - \eta \lambda_1 \in (-1, 1)$
        implies $\eta \lesssim \frac{1}{\lambda_1}$, roughly.
    \end{itemize}
    \item However, smaller values of $\eta$
    lead to slower convergence, as
    $1 - \eta \lambda_i$ is further from $0$
    when $\eta$ is smaller.
    \begin{itemize}
        \item The \emph{smallest} eigenvalue $\lambda_d$
        gives the slowest convergence here.
        For $\eta = \frac{1}{\lambda_1}$, the decay is:
        \[ \|w_t - w^*\|^2
        \leq \left(1 - \frac{\lambda_d}{\lambda_1}\right)^{2t}
        \|w_0 - w^*\|^2. \]
    \end{itemize}
\end{enumerate}

We can summarize this whole story like so.

\textbf{Definition.}
Let the \textbf{condition number} $\kappa$ of a matrix $H$
be $\kappa := \frac{\lambda_{\mathrm{max}}(H)}{\lambda_{\mathrm{min}}(H)}$
(the ratio of its largest eigenvalue to its smallest).

\begin{theorem}{Least-Squares Convergence Rate}
    Gradient descent for least-squares linear regression,
    with step size $\eta = \frac{1}{\lambda_{\mathrm{max}}(H)}$,
    converges to the global minimum
    at an exponential rate, described below.
    \[ \|w_t - w^*\|^2
    \leq \left(1 - \frac{1}{\kappa}\right)^{2t}
    \|w_0 - w^*\|^2. \]
\end{theorem}

This also explains why elliptically-shaped level sets
yield slower convergence:
they have a larger condition number.

\end{document}